\section{Inventaires Fongiques -- Complétude \&
Représentativité}\label{inventaires-fongiques-compluxe9tude-repruxe9sentativituxe9}

\href{https://www.r-project.org/}{\pandocbounded{\includesvg[keepaspectratio]{https://img.shields.io/badge/R-\%3E\%3D4.0-blue.svg}}}\\
\href{../scripts/Inventaires_completude_representativite.R}{\pandocbounded{\includesvg[keepaspectratio]{https://img.shields.io/badge/Script-Inventaires__completude__representativite.R-6f42c1.svg}}}\\
\href{../README.md}{\pandocbounded{\includesvg[keepaspectratio]{https://img.shields.io/badge/Domain-Mycologie\%20\%7C\%20\%C3\%89cologie-green.svg}}}\\
\hyperref[-quick-start]{\pandocbounded{\includesvg[keepaspectratio]{https://img.shields.io/badge/Status-Op\%C3\%A9rationnel-brightgreen.svg}}}

Pipeline R pour automatiser l'analyse de la \textbf{complétude
d'inventaires fongiques} et de leur \textbf{représentativité} (Taux
d'Espèces Exceptionnelles, \textbf{TEE} / Indice de représentativité,
\textbf{Ir}), avec sorties tabulaires et graphiques prêtes à exploiter.

\textbf{Projet} : Projet hôte / Inventaires fongiques\\
\textbf{Auteur} : Eddy Boite\\
\textbf{Version script} : 1.4 (selon en-tête du script)\\
\textbf{Date doc} : 22 Mars 2026\\
\textbf{Type de document} : Spécification fonctionnelle

\begin{center}\rule{0.5\linewidth}{0.5pt}\end{center}

\subsection{📖 À Propos}\label{uxe0-propos}

Ce document décrit l'utilisation du script :

\begin{itemize}
\tightlist
\item
  \texttt{scripts/Inventaires\_completude\_representativite.R}
\end{itemize}

Le script produit automatiquement, par site :

\begin{enumerate}
\def\labelenumi{\arabic{enumi}.}
\tightlist
\item
  Courbe temps-espèces
\item
  Ajustement linéaire et hyperbolique (si possible)
\item
  Estimation d'asymptote (richesse théorique)
\item
  Taux d'espèces exceptionnelles (TEE)
\item
  Indice de représentativité \(Ir = 1 - TEE\)
\item
  Fréquence des espèces par visite
\item
  Occupation spatiale (si \texttt{placette} disponible)
\item
  Analyse Factorielle des Correspondances (\textbf{AFC}, aussi appelée
  Correspondence Analysis, \textbf{CA}) simple (si package
  \texttt{vegan} disponible)
\item
  \textbf{Métriques de pertinence scientifique} : stabilité, robustesse,
  cohérence spatiotemporelle et score agrégé, avec visualisations
  dédiées
\end{enumerate}

En fin d'exécution, un \textbf{manifeste automatique} journalise l'état
de chaque fichier attendu (OK / MANQUANT / OPTIONNEL\_NON\_GENERE).

\subsection{Table des Matières}\label{table-des-matiuxe8res}

\begin{itemize}
\tightlist
\item
  \hyperref[-quick-start]{Quick Start}
\item
  \hyperref[-structure-du-projet]{Structure du projet}
\item
  \hyperref[-configurer-le-nom-du-projet-portabilitux5cux25C3ux5cux25A9]{Configurer
  le nom du projet (portabilité)}
\item
  \hyperref[-structure-attendue-des-donnux5cux25C3ux5cux25A9es]{Structure
  attendue des données}
\item
  \hyperref[-validation-csv-stricte-robustesse]{Validation CSV stricte
  (robustesse)}
\item
  \hyperref[-sorties-gux5cux25C3ux5cux25A9nux5cux25C3ux5cux25A9rux5cux25C3ux5cux25A9es]{Sorties
  générées}
\item
  \hyperref[-installation-et-dux5cux25C3ux5cux25A9pendances]{Installation
  et dépendances}
\item
  \hyperref[-configuration]{Configuration}
\item
  \hyperref[-interprux5cux25C3ux5cux25A9tation-des-indicateurs]{Interprétation
  des indicateurs}
\item
  \hyperref[-glossaire]{Glossaire}
\item
  \hyperref[rux5cux25C3ux5cux25A9fux5cux25C3ux5cux25A9rences]{Références}
\item
  \hyperref[-dux5cux25C3ux5cux25A9pannage-faq]{Dépannage (Foire Aux
  Questions, FAQ)}
\end{itemize}

\begin{center}\rule{0.5\linewidth}{0.5pt}\end{center}

\subsection{Structure du projet}\label{structure-du-projet}

Le script s'inscrit dans un projet R standard dont voici une structure
type :

\begin{verbatim}
<nom_du_projet>/
├── data/
│   ├── observations.csv                 # Données d'entrée ICR (à personnaliser)
├── docs/
│   ├── Readme_Analyse_Inventaires_Fongiques.md
│   ├── rapport_inventaires_fongiques.qmd          # Rapport Quarto
├── results/
│   ├── ICR_donnees_preparees.csv
│   ├── ICR_00_csv_conformite_report.csv
│   ├── ICR_00_csv_conformite_problems.csv
│   ├── ICR_resume_tous_sites.csv
│   ├── ICR_metrics_pertinence_tous_sites.csv
│   ├── ICR_comparaison_completude_sites.png
│   ├── ICR_comparaison_ir_sites.png
│   ├── ICR_metrics_pertinence_heatmap_sites.png   # Heatmap métriques multi-sites
│   ├── ICR_metrics_pertinence_score_sites.png     # Classement global pertinence
│   └── <site>/                          # Un sous-dossier par site
│       ├── ICR_01_courbe_temps_especes.csv
│       ├── ICR_02_tee_ir.csv
│       ├── ICR_03_frequence_especes.csv
│       ├── ICR_04_resume_site.csv
│       ├── ICR_05_modeles.csv           # Si modèles ajustés
│       ├── ICR_06_occupation_spatiale.csv  # Si placette disponible
│       ├── ICR_07_metrics_pertinence.csv
│       ├── ICR_01_richesse_par_visite_et_cumul.png
│       ├── ICR_02_courbe_temps_especes_hyperbole.png
│       ├── ICR_03_tee_ir.png
│       ├── ICR_04_histogramme_frequences.png
│       ├── ICR_05_temporel_vs_spatial.png
│       ├── ICR_06_CA_placettes_especes.png
│       └── ICR_07_metrics_pertinence_dashboard.png  # Dashboard métriques site
├── scripts/
│   └── Inventaires_completude_representativite.R  # ← Script principal (v1.4 Phase 4)
├── tests/
└── <nom_du_projet>.Rproj
\end{verbatim}

\begin{quote}
\textbf{Convention de nommage :} tous les fichiers générés par ce script
sont préfixés \texttt{ICR\_} pour les distinguer des sorties des autres
scripts du projet.
\end{quote}

\begin{center}\rule{0.5\linewidth}{0.5pt}\end{center}

\subsection{Configurer le nom du projet
(portabilité)}\label{configurer-le-nom-du-projet-portabilituxe9}

Ce document est conçu pour être \textbf{copié tel quel} dans un autre
dépôt.

\subsubsection{Ce qui dépend du nom du
projet}\label{ce-qui-duxe9pend-du-nom-du-projet}

Le script n'utilise \textbf{pas} le nom du dépôt dans son code métier.
Le « nom de projet » intervient surtout pour :

\begin{itemize}
\tightlist
\item
  l'affichage documentaire
  (\texttt{\textless{}nom\_du\_projet\textgreater{}}),
\item
  le nom du fichier RStudio
  (\texttt{\textless{}nom\_du\_projet\textgreater{}.Rproj}),
\item
  l'organisation de vos dossiers dans votre repo.
\end{itemize}

\subsubsection{Procédure recommandée (2
minutes)}\label{procuxe9dure-recommanduxe9e-2-minutes}

\begin{enumerate}
\def\labelenumi{\arabic{enumi}.}
\tightlist
\item
  Remplacez les placeholders
  \texttt{\textless{}nom\_du\_projet\textgreater{}} par votre vrai nom
  de dépôt (ex. \texttt{inventaires\_fongiques\_public}).
\item
  Si vous utilisez RStudio, renommez le fichier projet en conséquence :

  \begin{itemize}
  \tightlist
  \item
    \texttt{\textless{}nom\_du\_projet\textgreater{}.Rproj}
  \end{itemize}
\item
  Conservez la structure relative standard : \texttt{data/},
  \texttt{scripts/}, \texttt{results/}, \texttt{docs/}.
\end{enumerate}

\subsubsection{Configuration sans modifier le
script}\label{configuration-sans-modifier-le-script}

Le script supporte des variables d'environnement pour s'adapter à votre
projet :

\begin{itemize}
\tightlist
\item
  \texttt{INVENTAIRES\_INPUT\_FILE} : chemin d'entrée (défaut :
  \texttt{data/observations.csv})
\item
  \texttt{INVENTAIRES\_OUTPUT\_DIR} : dossier de sortie (défaut :
  \texttt{results/ICR})
\item
  \texttt{INVENTAIRES\_DATE\_FORMAT} : format de date (défaut :
  \texttt{\%Y-\%m-\%d})
\end{itemize}

Exemple d'usage portable :

\begin{itemize}
\tightlist
\item
  entrée personnalisée : \texttt{data/observations\_public.csv}
\item
  sorties personnalisées : \texttt{results/public\_release}
\item
  même script :
  \texttt{scripts/Inventaires\_completude\_representativite.R}
\end{itemize}

\subsubsection{Bonnes pratiques pour diffusion
publique}\label{bonnes-pratiques-pour-diffusion-publique}

\begin{itemize}
\tightlist
\item
  Évitez les chemins absolus locaux (\texttt{/Users/...}).
\item
  Gardez uniquement des chemins relatifs dans les docs et scripts.
\item
  Vérifiez que vos noms de dossiers ne révèlent pas d'information
  interne.
\end{itemize}

\begin{center}\rule{0.5\linewidth}{0.5pt}\end{center}

\subsection{Quick Start}\label{quick-start}

\subsubsection{1) Préparer les
données}\label{pruxe9parer-les-donnuxe9es}

Placer un fichier d'observations dans \texttt{data/} avec le nom
recommandé :

\begin{itemize}
\tightlist
\item
  \texttt{data/observations.csv}
\end{itemize}

\begin{quote}
Le script accepte aussi \texttt{observations.txt} ou
\texttt{observations.tsv} (\textbf{TSV} = Tab-Separated Values, valeurs
séparées par tabulation ; détection automatique du séparateur).
\end{quote}

\subsubsection{2) Lancer l'analyse}\label{lancer-lanalyse}

Depuis la racine du projet (quel que soit son nom) :

\begin{itemize}
\tightlist
\item
  \texttt{Rscript\ scripts/Inventaires\_completude\_representativite.R}
\end{itemize}

\subsubsection{3) Consulter les
résultats}\label{consulter-les-ruxe9sultats}

Les sorties sont produites dans :

\begin{itemize}
\tightlist
\item
  \texttt{results/}
\item
  puis un sous-dossier par site
  (\texttt{results/\textless{}site\_sanitize\textgreater{}/})
\end{itemize}

\begin{center}\rule{0.5\linewidth}{0.5pt}\end{center}

\subsection{Structure attendue des
données}\label{structure-attendue-des-donnuxe9es}

\subsubsection{Colonnes obligatoires}\label{colonnes-obligatoires}

\begin{itemize}
\tightlist
\item
  \texttt{site}
\item
  \texttt{date}
\item
  \texttt{visite\_id}
\item
  \texttt{espece}
\end{itemize}

\subsubsection{Colonne optionnelle}\label{colonne-optionnelle}

\begin{itemize}
\tightlist
\item
  \texttt{placette} (recommandée pour l'analyse spatiale et la CA)
\end{itemize}

\subsubsection{Format de date}\label{format-de-date}

Par défaut : \texttt{\%Y-\%m-\%d} (ex. \texttt{2026-03-18}).

\subsection{Validation CSV stricte
(robustesse)}\label{validation-csv-stricte-robustesse}

Le script applique un contrôle de conformité CSV avant les analyses :

\begin{enumerate}
\def\labelenumi{\arabic{enumi}.}
\tightlist
\item
  schéma de colonnes attendu,
\item
  contrôle des colonnes supplémentaires,
\item
  audit des problèmes de parsing remontés par
  \texttt{readr::problems()}.
\end{enumerate}

\subsubsection{Schéma attendu}\label{schuxe9ma-attendu}

\begin{itemize}
\tightlist
\item
  Colonnes requises (défaut) : \texttt{site}, \texttt{date},
  \texttt{visite\_id}, \texttt{espece}
\item
  Colonnes optionnelles (défaut) : \texttt{placette}
\end{itemize}

\subsubsection{Comportement en mode
strict}\label{comportement-en-mode-strict}

Par défaut, \texttt{csv\_strict\_mode\ =\ TRUE} (via
\texttt{INVENTAIRES\_CSV\_STRICT}).

Dans ce mode, l'exécution s'arrête si :

\begin{itemize}
\tightlist
\item
  une colonne obligatoire manque,
\item
  une colonne non autorisée est détectée,
\item
  un problème de parsing est présent.
\end{itemize}

\subsubsection{Rapports produits}\label{rapports-produits}

\begin{itemize}
\tightlist
\item
  \texttt{results/ICR/ICR\_00\_csv\_conformite\_report.csv} (toujours)
\item
  \texttt{results/ICR/ICR\_00\_csv\_conformite\_problems.csv} (si
  anomalies de parsing)
\end{itemize}

\subsubsection{Paramètres et variables
d'environnement}\label{paramuxe8tres-et-variables-denvironnement}

\begin{itemize}
\tightlist
\item
  \texttt{CONFIG\$csv\_strict\_mode} ↔ \texttt{INVENTAIRES\_CSV\_STRICT}
  (défaut \texttt{TRUE})
\item
  \texttt{CONFIG\$csv\_allow\_extra\_cols} ↔
  \texttt{INVENTAIRES\_CSV\_ALLOW\_EXTRA\_COLS} (défaut \texttt{FALSE})
\item
  \texttt{CONFIG\$csv\_required\_cols} (défaut :
  \texttt{site,date,visite\_id,espece})
\item
  \texttt{CONFIG\$csv\_optional\_cols} (défaut : \texttt{placette})
\end{itemize}

\subsubsection{Définition d'une visite distincte
(important)}\label{duxe9finition-dune-visite-distincte-important}

Le script considère une visite unique via la clé :

\begin{itemize}
\tightlist
\item
  \texttt{date\ +\ visite\_id} (et \texttt{site} lorsque plusieurs sites
  sont analysés)
\end{itemize}

Cela permet de gérer les jeux de données où un même \texttt{visite\_id}
peut être réutilisé à des dates différentes.

\subsubsection{Exemple minimal}\label{exemple-minimal}

{\def\LTcaptype{none} % do not increment counter
\begin{longtable}[]{@{}lllll@{}}
\toprule\noalign{}
site & date & visite\_id & espece & placette \\
\midrule\noalign{}
\endhead
\bottomrule\noalign{}
\endlastfoot
Site\_A & 2026-03-01 & V1 & Amanita muscaria & P1 \\
Site\_A & 2026-03-01 & V1 & Boletus edulis & P2 \\
Site\_A & 2026-03-15 & V2 & Amanita muscaria & P1 \\
Site\_B & 2026-03-08 & V1 & Cantharellus cibarius & P3 \\
\end{longtable}
}

\begin{center}\rule{0.5\linewidth}{0.5pt}\end{center}

\subsection{Sorties générées}\label{sorties-guxe9nuxe9ruxe9es}

\subsubsection{Vue d'ensemble du pipeline (9 calculs →
fichiers)}\label{vue-densemble-du-pipeline-9-calculs-fichiers}

\begin{verbatim}
1) Préparation/validation des données
    -> Global CSV : ICR_donnees_preparees.csv
2) Richesse observée et cumulée (courbe temps-espèces)
    -> Site CSV   : <site>/ICR_01_courbe_temps_especes.csv
    -> Site PNG   : <site>/ICR_01_richesse_par_visite_et_cumul.png
3) Ajustements linéaire/hyperbolique (si conditions remplies)
    -> Site CSV   : <site>/ICR_05_modeles.csv
    -> Site PNG   : <site>/ICR_02_courbe_temps_especes_hyperbole.png
4) Asymptote hyperbolique et ratio de complétude
    -> Site CSV   : <site>/ICR_04_resume_site.csv
    -> Global CSV : ICR_resume_tous_sites.csv
    -> Global PNG : ICR_comparaison_completude_sites.png
5) TEE (taux d'espèces exceptionnelles) et indice Ir
    -> Site CSV   : <site>/ICR_02_tee_ir.csv
    -> Site PNG   : <site>/ICR_03_tee_ir.png
    -> Global PNG : ICR_comparaison_ir_sites.png
6) Fréquences d'occurrence des espèces
    -> Site CSV   : <site>/ICR_03_frequence_especes.csv
    -> Site PNG   : <site>/ICR_04_histogramme_frequences.png
7) Occupation spatiale (si placette disponible)
    -> Site CSV   : <site>/ICR_06_occupation_spatiale.csv
    -> Site PNG   : <site>/ICR_05_temporel_vs_spatial.png
8) CA/AFC placettes-espèces (si vegan et données adaptées)
    -> Site PNG   : <site>/ICR_06_CA_placettes_especes.png
9) Métriques de pertinence scientifique (stabilité, robustesse, cohérence)
    -> Site CSV   : <site>/ICR_07_metrics_pertinence.csv
    -> Global CSV : ICR_metrics_pertinence_tous_sites.csv
    -> Site PNG   : <site>/ICR_07_metrics_pertinence_dashboard.png
    -> Global PNG : ICR_metrics_pertinence_heatmap_sites.png
    -> Global PNG : ICR_metrics_pertinence_score_sites.png
\end{verbatim}

\begin{quote}
En fin d'exécution, un \textbf{manifeste automatique} liste l'état de
chaque fichier attendu : OK / MANQUANT / OPTIONNEL\_NON\_GENERE.
\end{quote}

\begin{center}\rule{0.5\linewidth}{0.5pt}\end{center}

\subsubsection{\texorpdfstring{Sorties globales
(\texttt{results/})}{Sorties globales (results/)}}\label{sorties-globales-results}

{\def\LTcaptype{none} % do not increment counter
\begin{longtable}[]{@{}
  >{\raggedright\arraybackslash}p{(\linewidth - 4\tabcolsep) * \real{0.3333}}
  >{\raggedright\arraybackslash}p{(\linewidth - 4\tabcolsep) * \real{0.3333}}
  >{\raggedright\arraybackslash}p{(\linewidth - 4\tabcolsep) * \real{0.3333}}@{}}
\toprule\noalign{}
\begin{minipage}[b]{\linewidth}\raggedright
Fichier
\end{minipage} & \begin{minipage}[b]{\linewidth}\raggedright
Type
\end{minipage} & \begin{minipage}[b]{\linewidth}\raggedright
Description
\end{minipage} \\
\midrule\noalign{}
\endhead
\bottomrule\noalign{}
\endlastfoot
\texttt{ICR\_00\_csv\_conformite\_report.csv} & CSV & Rapport de
conformité CSV (schéma + parsing) avant traitement. \\
\texttt{ICR\_00\_csv\_conformite\_problems.csv} & CSV & Détail des
anomalies de parsing détectées (si présentes). \\
\texttt{ICR\_donnees\_preparees.csv} & CSV & Données nettoyées et
standardisées utilisées pour l'analyse. \\
\texttt{ICR\_resume\_tous\_sites.csv} & CSV & Synthèse finale
multi-sites (complétude, représentativité, etc.). \\
\texttt{ICR\_metrics\_pertinence\_tous\_sites.csv} & CSV & Métriques de
pertinence scientifique agrégées pour tous les sites. \\
\texttt{ICR\_comparaison\_completude\_sites.png} & Graphique &
Comparaison du ratio de complétude entre sites. \\
\texttt{ICR\_comparaison\_ir\_sites.png} & Graphique & Comparaison de
l'indice \(Ir\) final entre sites. \\
\texttt{ICR\_metrics\_pertinence\_heatmap\_sites.png} & Graphique &
Heatmap comparative des 7 métriques de pertinence pour tous les
sites. \\
\texttt{ICR\_metrics\_pertinence\_score\_sites.png} & Graphique &
Classement global des sites par score de pertinence scientifique
agrégé. \\
\end{longtable}
}

\subsubsection{\texorpdfstring{Sorties par site
(\texttt{results/\textless{}site\textgreater{}/})}{Sorties par site (results/\textless site\textgreater/)}}\label{sorties-par-site-resultssite}

{\def\LTcaptype{none} % do not increment counter
\begin{longtable}[]{@{}
  >{\raggedright\arraybackslash}p{(\linewidth - 6\tabcolsep) * \real{0.2500}}
  >{\raggedright\arraybackslash}p{(\linewidth - 6\tabcolsep) * \real{0.2500}}
  >{\raggedright\arraybackslash}p{(\linewidth - 6\tabcolsep) * \real{0.2500}}
  >{\raggedright\arraybackslash}p{(\linewidth - 6\tabcolsep) * \real{0.2500}}@{}}
\toprule\noalign{}
\begin{minipage}[b]{\linewidth}\raggedright
Fichier
\end{minipage} & \begin{minipage}[b]{\linewidth}\raggedright
Type
\end{minipage} & \begin{minipage}[b]{\linewidth}\raggedright
Condition
\end{minipage} & \begin{minipage}[b]{\linewidth}\raggedright
Description
\end{minipage} \\
\midrule\noalign{}
\endhead
\bottomrule\noalign{}
\endlastfoot
\texttt{ICR\_01\_courbe\_temps\_especes.csv} & CSV & Toujours & Données
de richesse observée et cumulée par visite. \\
\texttt{ICR\_02\_tee\_ir.csv} & CSV & Toujours & Évolution de TEE et de
\(Ir\) au fil des visites. \\
\texttt{ICR\_03\_frequence\_especes.csv} & CSV & Toujours & Fréquences
d'occurrence des espèces par visite. \\
\texttt{ICR\_04\_resume\_site.csv} & CSV & Toujours & Résumé synthétique
du site analysé. \\
\texttt{ICR\_05\_modeles.csv} & CSV & Si modèles ajustés & Paramètres et
qualité d'ajustement des modèles. \\
\texttt{ICR\_06\_occupation\_spatiale.csv} & CSV & Si \texttt{placette}
exploitable & Occupation spatiale des espèces par placette. \\
\texttt{ICR\_01\_richesse\_par\_visite\_et\_cumul.png} & Graphique &
Toujours & Richesse par visite + richesse cumulée. \\
\texttt{ICR\_02\_courbe\_temps\_especes\_hyperbole.png} & Graphique &
Toujours (courbe) / hyperbole si modèle dispo & Courbe temps-espèces
avec ajustement hyperbolique si possible. \\
\texttt{ICR\_03\_tee\_ir.png} & Graphique & Toujours & Lecture conjointe
espèces uniques, total et \(Ir\). \\
\texttt{ICR\_04\_histogramme\_frequences.png} & Graphique & Toujours &
Distribution des fréquences d'observation des espèces. \\
\texttt{ICR\_05\_temporel\_vs\_spatial.png} & Graphique & Si spatial
disponible & Relation fréquence temporelle vs occupation spatiale. \\
\texttt{ICR\_06\_CA\_placettes\_especes.png} & Graphique & Si
\texttt{vegan} + données adaptées & Ordination AFC/CA
placettes-espèces. \\
\texttt{ICR\_07\_metrics\_pertinence.csv} & CSV & Toujours & Métriques
de pertinence scientifique (15 colonnes : Ir, complétude, R²,
score\ldots). \\
\texttt{ICR\_07\_metrics\_pertinence\_dashboard.png} & Graphique &
Toujours & Dashboard des 7 métriques de pertinence scientifique du site
(stabilité, robustesse, cohérence). \\
\end{longtable}
}

\begin{center}\rule{0.5\linewidth}{0.5pt}\end{center}

\subsection{Installation et
dépendances}\label{installation-et-duxe9pendances}

\subsubsection{Packages requis (base)}\label{packages-requis-base}

{\def\LTcaptype{none} % do not increment counter
\begin{longtable}[]{@{}
  >{\raggedright\arraybackslash}p{(\linewidth - 2\tabcolsep) * \real{0.5000}}
  >{\raggedright\arraybackslash}p{(\linewidth - 2\tabcolsep) * \real{0.5000}}@{}}
\toprule\noalign{}
\begin{minipage}[b]{\linewidth}\raggedright
Package
\end{minipage} & \begin{minipage}[b]{\linewidth}\raggedright
Rôle principal
\end{minipage} \\
\midrule\noalign{}
\endhead
\bottomrule\noalign{}
\endlastfoot
\texttt{dplyr} & Manipulation de données tabulaires (filtre, jointures,
agrégations). \\
\texttt{tidyr} & Restructuration des données (format long/large). \\
\texttt{ggplot2} & Génération des graphiques. \\
\texttt{purrr} & Itérations fonctionnelles sur les sites/listes. \\
\texttt{readr} & Lecture/écriture rapide des fichiers texte/CSV. \\
\texttt{stringr} & Traitements de chaînes de caractères. \\
\texttt{forcats} & Gestion des facteurs et de leurs niveaux. \\
\texttt{tibble} & Structures tabulaires modernes. \\
\texttt{scales} & Mise en forme des axes et échelles graphiques. \\
\texttt{gridExtra} & Assemblage de graphiques multi-panneaux (notamment
\texttt{ICR\_03\_tee\_ir.png}). \\
\end{longtable}
}

\subsubsection{Packages optionnels}\label{packages-optionnels}

{\def\LTcaptype{none} % do not increment counter
\begin{longtable}[]{@{}
  >{\raggedright\arraybackslash}p{(\linewidth - 2\tabcolsep) * \real{0.5000}}
  >{\raggedright\arraybackslash}p{(\linewidth - 2\tabcolsep) * \real{0.5000}}@{}}
\toprule\noalign{}
\begin{minipage}[b]{\linewidth}\raggedright
Package
\end{minipage} & \begin{minipage}[b]{\linewidth}\raggedright
Utilisation
\end{minipage} \\
\midrule\noalign{}
\endhead
\bottomrule\noalign{}
\endlastfoot
\texttt{minpack.lm} & Ajustement du modèle hyperbolique (estimation
d'asymptote). \\
\texttt{vegan} & Ordination AFC/CA (analyse placettes--espèces). \\
\end{longtable}
}

Le script vérifie les packages requis au démarrage et affiche une
commande d'installation explicite si certains sont manquants.

\begin{center}\rule{0.5\linewidth}{0.5pt}\end{center}

\subsection{Configuration}\label{configuration}

La configuration est centralisée dans l'objet \texttt{CONFIG} du script.

{\def\LTcaptype{none} % do not increment counter
\begin{longtable}[]{@{}
  >{\raggedright\arraybackslash}p{(\linewidth - 4\tabcolsep) * \real{0.3333}}
  >{\raggedright\arraybackslash}p{(\linewidth - 4\tabcolsep) * \real{0.3333}}
  >{\raggedright\arraybackslash}p{(\linewidth - 4\tabcolsep) * \real{0.3333}}@{}}
\toprule\noalign{}
\begin{minipage}[b]{\linewidth}\raggedright
Paramètre
\end{minipage} & \begin{minipage}[b]{\linewidth}\raggedright
Défaut
\end{minipage} & \begin{minipage}[b]{\linewidth}\raggedright
Description
\end{minipage} \\
\midrule\noalign{}
\endhead
\bottomrule\noalign{}
\endlastfoot
\texttt{input\_file} & \texttt{data/observations.csv} & Chemin du
fichier d'entrée des observations. \\
\texttt{output\_dir} & \texttt{results/ICR} & Dossier de sortie pour les
CSV et graphiques. \\
\texttt{date\_format} & \texttt{\%Y-\%m-\%d} & Format attendu pour
parser la colonne \texttt{date}. \\
\texttt{csv\_strict\_mode} & \texttt{TRUE} & Stoppe l'exécution si le
CSV n'est pas conforme. \\
\texttt{csv\_allow\_extra\_cols} & \texttt{FALSE} & Autorise/refuse les
colonnes supplémentaires non prévues. \\
\texttt{csv\_required\_cols} &
\texttt{c("site","date","visite\_id","espece")} & Colonnes obligatoires
du schéma CSV. \\
\texttt{csv\_optional\_cols} & \texttt{c("placette")} & Colonnes
optionnelles autorisées. \\
\texttt{make\_ca} & \texttt{TRUE} & Active l'ordination AFC/CA si les
conditions sont réunies. \\
\texttt{min\_visits\_for\_model} & \texttt{5} & Nombre minimal de
visites pour ajuster les modèles de complétude. \\
\texttt{width} & \texttt{10} & Largeur (en pouces) des graphiques
exportés. \\
\texttt{height} & \texttt{6} & Hauteur (en pouces) des graphiques
exportés. \\
\texttt{dpi} & \texttt{300} & Résolution d'export des graphiques (dots
per inch). \\
\end{longtable}
}

Le script accepte aussi des surcharges via variables d'environnement :

\begin{itemize}
\tightlist
\item
  \texttt{INVENTAIRES\_INPUT\_FILE}
\item
  \texttt{INVENTAIRES\_OUTPUT\_DIR}
\item
  \texttt{INVENTAIRES\_DATE\_FORMAT}
\item
  \texttt{INVENTAIRES\_CSV\_STRICT}
\item
  \texttt{INVENTAIRES\_CSV\_ALLOW\_EXTRA\_COLS}
\end{itemize}

Si besoin, adapter ces valeurs directement dans :

\begin{itemize}
\tightlist
\item
  \texttt{scripts/Inventaires\_completude\_representativite.R}
\end{itemize}

\begin{center}\rule{0.5\linewidth}{0.5pt}\end{center}

\subsection{Interprétation des
indicateurs}\label{interpruxe9tation-des-indicateurs}

\subsubsection{TEE (Taux d'Espèces
Exceptionnelles)}\label{tee-taux-despuxe8ces-exceptionnelles}

\begin{itemize}
\tightlist
\item
  Définit la proportion d'espèces observées une seule fois.
\item
  Plus TEE est élevé, plus l'inventaire peut être incomplet.
\end{itemize}

Lecture pratique :

\begin{itemize}
\tightlist
\item
  \textbf{TEE élevé en fin de suivi} : beaucoup d'espèces n'ont été vues
  qu'une fois, ce qui suggère un inventaire encore en phase
  d'accumulation.
\item
  \textbf{TEE qui diminue avec les visites} : signe positif, les espèces
  ``uniques'' deviennent progressivement récurrentes.
\item
  \textbf{TEE stable et bas} : communauté mieux couverte, effort
  d'échantillonnage proche d'un plateau.
\end{itemize}

Repères empiriques (à contextualiser selon habitat/saison/protocole) :

\begin{itemize}
\tightlist
\item
  TEE \textgreater{} 0.40 : inventaire souvent \textbf{incomplet}
\item
  0.20 ≤ TEE ≤ 0.40 : inventaire \textbf{intermédiaire}
\item
  TEE \textless{} 0.20 : inventaire généralement \textbf{bien consolidé}
\end{itemize}

\subsubsection{Ir (Indice de
représentativité)}\label{ir-indice-de-repruxe9sentativituxe9}

\[Ir = 1 - TEE\]

\begin{itemize}
\tightlist
\item
  Plus \(Ir\) est proche de 1, meilleure est la représentativité de
  l'inventaire.
\end{itemize}

Lecture pratique :

\begin{itemize}
\tightlist
\item
  \textbf{Ir en hausse au fil des visites} : le protocole capture mieux
  la diversité réelle du site.
\item
  \textbf{Ir qui plafonne tôt} : effort potentiellement suffisant, sauf
  si période écologique incomplète.
\item
  \textbf{Ir faible et stagnant} : effort, calendrier ou stratification
  spatiale à revoir.
\end{itemize}

Repères empiriques (indicatifs) :

\begin{itemize}
\tightlist
\item
  Ir \textless{} 0.60 : représentativité faible
\item
  0.60 ≤ Ir \textless{} 0.80 : représentativité moyenne
\item
  Ir ≥ 0.80 : représentativité bonne à élevée
\end{itemize}

\subsubsection{Complétude observée /
asymptote}\label{compluxe9tude-observuxe9e-asymptote}

Le script estime une asymptote de richesse (modèle hyperbolique, si
disponible) puis calcule un ratio de complétude :

\[\\text{Complétude} = \\frac{S\_{obs}}{S\_{asymptote}}\]

avec :

\begin{itemize}
\tightlist
\item
  \(S\_{obs}\) : richesse observée finale
\item
  \(S\_{asymptote}\) : richesse théorique estimée
\end{itemize}

Interprétation :

\begin{itemize}
\tightlist
\item
  \textbf{Complétude proche de 1} : l'inventaire approche la richesse
  théorique attendue.
\item
  \textbf{Complétude intermédiaire} : des espèces restent probablement à
  détecter.
\item
  \textbf{Complétude faible} : effort insuffisant ou forte hétérogénéité
  écologique.
\end{itemize}

Repères utiles :

\begin{itemize}
\tightlist
\item
  Complétude \textless{} 0.70 : insuffisante pour conclure sur
  l'exhaustivité
\item
  0.70--0.90 : progression correcte, inventaire partiel
\end{itemize}

\begin{quote}
0.90 : inventaire très avancé
\end{quote}

\begin{center}\rule{0.5\linewidth}{0.5pt}\end{center}

\subsubsection{Lecture détaillée des diagrammes (le plus
important)}\label{lecture-duxe9tailluxe9e-des-diagrammes-le-plus-important}

\paragraph{\texorpdfstring{\texttt{ICR\_01\_richesse\_par\_visite\_et\_cumul.png}}{ICR\_01\_richesse\_par\_visite\_et\_cumul.png}}\label{icr_01_richesse_par_visite_et_cumul.png}

Ce graphique combine :

\begin{itemize}
\tightlist
\item
  \textbf{barres} = richesse observée par visite
  (\texttt{nb\_especes\_trouvees})
\item
  \textbf{ligne} = richesse cumulée (\texttt{nb\_cumule})
\end{itemize}

Comment lire :

\begin{itemize}
\tightlist
\item
  \textbf{Barres encore hautes + cumul qui monte vite} : inventaire en
  phase active de découverte.
\item
  \textbf{Barres qui baissent + cumul qui se tasse} : entrée dans une
  phase de saturation.
\item
  \textbf{Dents de scie marquées} : forte dépendance
  météo/saison/observateur.
\end{itemize}

Signal d'alerte :

\begin{itemize}
\tightlist
\item
  Si les dernières visites apportent encore beaucoup de nouvelles
  espèces, l'arrêt de l'inventaire peut être prématuré.
\end{itemize}

\paragraph{\texorpdfstring{\texttt{ICR\_02\_courbe\_temps\_especes\_hyperbole.png}}{ICR\_02\_courbe\_temps\_especes\_hyperbole.png}}\label{icr_02_courbe_temps_especes_hyperbole.png}

Ce graphique superpose les observations à un ajustement hyperbolique (si
\texttt{minpack.lm} disponible) et une asymptote.

Comment lire :

\begin{itemize}
\tightlist
\item
  \textbf{Courbe observée proche de la courbe ajustée} : modèle
  cohérent.
\item
  \textbf{Écart important} : structure de données non compatible avec ce
  modèle (saisonnalité forte, ruptures d'effort, etc.).
\item
  \textbf{Asymptote nettement au-dessus de la richesse actuelle} : marge
  de découverte encore importante.
\end{itemize}

Point clé :

\begin{itemize}
\tightlist
\item
  Ne pas interpréter l'asymptote comme une ``vérité absolue'', mais
  comme un \textbf{repère quantitatif} de complétude.
\end{itemize}

\paragraph{\texorpdfstring{\texttt{ICR\_03\_tee\_ir.png}}{ICR\_03\_tee\_ir.png}}\label{icr_03_tee_ir.png}

Ce graphique présente \textbf{deux panneaux complémentaires} :

\begin{itemize}
\tightlist
\item
  panneau haut : nombre d'espèces observées une seule fois
  (\texttt{nb\_esp\_1fois}),
\item
  panneau bas : indice de représentativité \(Ir\) sur l'échelle
  \(\[0,1\]\) avec seuils visuels.
\end{itemize}

Comment lire :

\begin{itemize}
\tightlist
\item
  \textbf{Espèces uniques qui ralentissent} : signal de consolidation de
  l'inventaire.
\item
  \textbf{Ir qui augmente puis se stabilise} : représentativité qui
  converge.
\item
  \textbf{Ir sous 0.60} : représentativité faible ; \textbf{0.60--0.80}
  : moyenne ; \textbf{≥ 0.80} : bonne.
\end{itemize}

Point de vigilance :

\begin{itemize}
\tightlist
\item
  Les seuils sont des repères empiriques : toujours interpréter avec le
  contexte (saison, effort, habitat).
\end{itemize}

\paragraph{\texorpdfstring{\texttt{ICR\_04\_histogramme\_frequences.png}}{ICR\_04\_histogramme\_frequences.png}}\label{icr_04_histogramme_frequences.png}

Histogramme du nombre d'espèces selon leur fréquence d'occurrence
(nombre de visites où elles sont observées).

Comment lire :

\begin{itemize}
\tightlist
\item
  \textbf{Pic sur les faibles fréquences (1--2 visites)} : communauté
  riche en espèces rares/occasionnelles.
\item
  \textbf{Répartition plus équilibrée} : noyau d'espèces régulières bien
  documenté.
\item
  \textbf{Queue vers les fortes fréquences} : présence d'espèces
  constantes, potentiels bons indicateurs écologiques.
\end{itemize}

Usages :

\begin{itemize}
\tightlist
\item
  Ajuster l'effort d'échantillonnage selon l'objectif (détection rareté
  vs suivi d'espèces communes).
\end{itemize}

\paragraph{\texorpdfstring{\texttt{ICR\_05\_temporel\_vs\_spatial.png}}{ICR\_05\_temporel\_vs\_spatial.png}}\label{icr_05_temporel_vs_spatial.png}

Nuage de points reliant :

\begin{itemize}
\tightlist
\item
  \(x\) = proportion de visites où l'espèce est vue
  (\texttt{prop\_visites})
\item
  \(y\) = proportion de placettes occupées (\texttt{prop\_placettes})
\end{itemize}

Lecture écologique :

\begin{itemize}
\tightlist
\item
  \textbf{Haut-droite} : espèces fréquentes et largement réparties
  (généralistes/communes).
\item
  \textbf{Bas-gauche} : espèces rares et localisées.
\item
  \textbf{Bas-droite} : espèces temporellement fréquentes mais
  spatialement confinées (micro-habitats stables).
\item
  \textbf{Haut-gauche} : espèces spatialement larges mais irrégulières
  dans le temps (phénologie/pulses).
\end{itemize}

Ce diagramme est très utile pour distinguer :

\begin{itemize}
\tightlist
\item
  rareté vraie,
\item
  rareté d'observation,
\item
  et hétérogénéité spatiale.
\end{itemize}

\paragraph{\texorpdfstring{\texttt{ICR\_06\_CA\_placettes\_especes.png}
(si
disponible)}{ICR\_06\_CA\_placettes\_especes.png (si disponible)}}\label{icr_06_ca_placettes_especes.png-si-disponible}

Ordination CA/AFC pour explorer les associations placettes--espèces.

Comment lire :

\begin{itemize}
\tightlist
\item
  \textbf{Points proches} : profils floristiques/mycologiques
  similaires.
\item
  \textbf{Points éloignés} : composition différente.
\item
  \textbf{Groupes compacts} : unités écologiques homogènes.
\item
  \textbf{Points isolés} : placettes atypiques ou espèces spécialisées.
\end{itemize}

Précaution d'interprétation :

\begin{itemize}
\tightlist
\item
  Les axes de CA sont relatifs au jeu de données analysé ; comparer des
  CA entre jeux différents demande prudence.
\end{itemize}

\begin{center}\rule{0.5\linewidth}{0.5pt}\end{center}

\subsubsection{Lecture des comparaisons
multi-sites}\label{lecture-des-comparaisons-multi-sites}

\paragraph{\texorpdfstring{\texttt{ICR\_comparaison\_completude\_sites.png}}{ICR\_comparaison\_completude\_sites.png}}\label{icr_comparaison_completude_sites.png}

Classe les sites par ratio de complétude \(S\_{obs}/S\_{asymptote}\).

\begin{itemize}
\tightlist
\item
  Permet d'identifier rapidement les sites à renforcer en effort
  d'inventaire.
\item
  Un site bas en complétude n'est pas ``mauvais'' : il peut être plus
  hétérogène ou plus difficile à prospecter.
\end{itemize}

\paragraph{\texorpdfstring{\texttt{ICR\_comparaison\_ir\_sites.png}}{ICR\_comparaison\_ir\_sites.png}}\label{icr_comparaison_ir_sites.png}

Classe les sites par \(Ir\) final.

\begin{itemize}
\tightlist
\item
  Très utile pour prioriser les suivis : sites à \(Ir\) faible =
  priorité de revisite.
\item
  Croiser avec la complétude est recommandé pour éviter les conclusions
  hâtives.
\end{itemize}

\begin{center}\rule{0.5\linewidth}{0.5pt}\end{center}

\paragraph{\texorpdfstring{\texttt{ICR\_metrics\_pertinence\_heatmap\_sites.png}}{ICR\_metrics\_pertinence\_heatmap\_sites.png}}\label{icr_metrics_pertinence_heatmap_sites.png}

Heatmap comparative (7 métriques × tous les sites), couleurs du rouge
(faible) au vert (bon).

\begin{itemize}
\tightlist
\item
  Chaque cellule affiche la valeur normalisée (0--1) de l'indicateur
  pour le site.
\item
  Permet d'identifier en un coup d'œil les métriques défaillantes et les
  sites les plus solides.
\item
  7 métriques visualisées : \(Ir\) final, complétude, stabilité finale,
  robustesse modèle R², cohérence tempo-spatiale Spearman, concordance
  tempo-spatiale, score de pertinence global.
\end{itemize}

\paragraph{\texorpdfstring{\texttt{ICR\_metrics\_pertinence\_score\_sites.png}}{ICR\_metrics\_pertinence\_score\_sites.png}}\label{icr_metrics_pertinence_score_sites.png}

Classement horizontal des sites par \textbf{score de pertinence agrégé}
(pondération : Ir 35 \%, complétude 35 \%, stabilité 20 \%, R² 10 \%).

\begin{itemize}
\tightlist
\item
  Zones de fond : rouge (\textless{} 60 \% = faible), jaune (60--80 \% =
  intermédiaire), vert (\textgreater{} 80 \% = bon).
\item
  Synthèse exécutive idéale pour prioriser les sites nécessitant un
  effort supplémentaire.
\end{itemize}

\paragraph{\texorpdfstring{\texttt{ICR\_07\_metrics\_pertinence\_dashboard.png}
(par
site)}{ICR\_07\_metrics\_pertinence\_dashboard.png (par site)}}\label{icr_07_metrics_pertinence_dashboard.png-par-site}

Dashboard panoramique par site : barres horizontales pour chacun des 7
indicateurs, normalisés de 0 à 1.

\begin{itemize}
\tightlist
\item
  Lignes de seuil pointillées à 0,60 (acceptable) et 0,80 (bon).
\item
  Lecture : tout indicateur sous 0,60 motive un complément d'effort ou
  une révision du protocole.
\end{itemize}

\begin{center}\rule{0.5\linewidth}{0.5pt}\end{center}

\subsubsection{Bonnes pratiques
d'interprétation}\label{bonnes-pratiques-dinterpruxe9tation}

\begin{itemize}
\tightlist
\item
  Interpréter \textbf{ensemble} : courbe cumulée + TEE/Ir + fréquences +
  spatial.
\item
  Toujours contextualiser avec : saison, pression d'échantillonnage,
  météo, observateurs.
\item
  Préférer les tendances sur plusieurs visites plutôt qu'une lecture
  ponctuelle.
\item
  Pour décisions de gestion, combiner ces indicateurs avec expertise
  terrain et connaissance taxonomique.
\end{itemize}

\begin{center}\rule{0.5\linewidth}{0.5pt}\end{center}

\subsection{Dépannage (FAQ)}\label{duxe9pannage-faq}

\subsubsection{Le manifeste de fin d'exécution signale des fichiers
MANQUANT}\label{le-manifeste-de-fin-dexuxe9cution-signale-des-fichiers-manquant}

Le manifeste journalise l'état de chaque fichier attendu à la fin de
l'analyse (statuts : OK, MANQUANT, OPTIONNEL\_NON\_GENERE). Un statut
MANQUANT indique qu'un fichier obligatoire n'a pas été produit :

\begin{itemize}
\tightlist
\item
  Vérifier les logs d'exécution pour l'erreur associée.
\item
  Les fichiers \texttt{OPTIONNEL\_NON\_GENERE} sont normaux si les
  conditions d'activation ne sont pas réunies (ex. \texttt{placette}
  absent, \texttt{vegan} non installé, ou \texttt{make\_ca\ =\ FALSE}).
\end{itemize}

\subsubsection{Les graphiques de métriques ne sont pas
générés}\label{les-graphiques-de-muxe9triques-ne-sont-pas-guxe9nuxe9ruxe9s}

Les 3 graphiques de métriques de pertinence nécessitent que
\texttt{ICR\_07\_metrics\_pertinence.csv} ait été produit avec des
données valides. Vérifier :

\begin{itemize}
\tightlist
\item
  Que le site comporte au moins 2 visites.
\item
  Que les colonnes \texttt{Ir\_final} et \texttt{completude} sont
  présentes dans le résumé du site.
\end{itemize}

\subsubsection{Le script ne trouve pas mon fichier
d'entrée}\label{le-script-ne-trouve-pas-mon-fichier-dentruxe9e}

Vérifier qu'un fichier existe dans \texttt{data/} sous l'un des noms
suivants :

\begin{itemize}
\tightlist
\item
  \texttt{observations.csv}
\item
  \texttt{observations.txt}
\item
  \texttt{observations.tsv}
\end{itemize}

et que les colonnes obligatoires sont présentes.

\subsubsection{Le script s'arrête avec ``CSV non
conforme''}\label{le-script-sarruxeate-avec-csv-non-conforme}

Consulter en priorité :

\begin{itemize}
\tightlist
\item
  \texttt{results/ICR/ICR\_00\_csv\_conformite\_report.csv}
\item
  \texttt{results/ICR/ICR\_00\_csv\_conformite\_problems.csv} (si
  présent)
\end{itemize}

Causes typiques :

\begin{itemize}
\tightlist
\item
  colonne obligatoire absente (\texttt{site}, \texttt{date},
  \texttt{visite\_id}, \texttt{espece}),
\item
  colonne supplémentaire non autorisée
  (\texttt{csv\_allow\_extra\_cols\ =\ FALSE}),
\item
  ligne mal formée (parse failure).
\end{itemize}

Pour assouplir temporairement :

\begin{itemize}
\tightlist
\item
  \texttt{INVENTAIRES\_CSV\_STRICT=FALSE}
\item
  \texttt{INVENTAIRES\_CSV\_ALLOW\_EXTRA\_COLS=TRUE}
\end{itemize}

\subsubsection{J'ai une erreur sur les
dates}\label{jai-une-erreur-sur-les-dates}

Contrôler le format de la colonne \texttt{date} et l'option
\texttt{CONFIG\$date\_format}.

\subsubsection{Le modèle hyperbolique n'est pas
calculé}\label{le-moduxe8le-hyperbolique-nest-pas-calculuxe9}

Installer \texttt{minpack.lm} (optionnel). Sans ce package, l'analyse
continue mais sans asymptote hyperbolique.

\subsubsection{La CA/AFC n'apparaît
pas}\label{la-caafc-napparauxeet-pas}

Nécessite :

\begin{itemize}
\tightlist
\item
  package \texttt{vegan}
\item
  colonne \texttt{placette} utilisable
\item
  matrice placette-espèce suffisamment informative
\end{itemize}

\begin{center}\rule{0.5\linewidth}{0.5pt}\end{center}

\subsection{🧾 Glossaire}\label{glossaire}

{\def\LTcaptype{none} % do not increment counter
\begin{longtable}[]{@{}
  >{\raggedright\arraybackslash}p{(\linewidth - 2\tabcolsep) * \real{0.5000}}
  >{\raggedright\arraybackslash}p{(\linewidth - 2\tabcolsep) * \real{0.5000}}@{}}
\toprule\noalign{}
\begin{minipage}[b]{\linewidth}\raggedright
Acronyme
\end{minipage} & \begin{minipage}[b]{\linewidth}\raggedright
Définition
\end{minipage} \\
\midrule\noalign{}
\endhead
\bottomrule\noalign{}
\endlastfoot
\textbf{TEE} & \textbf{Taux d'Espèces Exceptionnelles} : proportion
d'espèces observées une seule fois dans l'inventaire. \\
\textbf{Ir} & \textbf{Indice de représentativité} : indicateur
synthétique défini par \(Ir = 1 - TEE\). \\
\textbf{Robustesse} & Qualité d'ajustement du modèle d'asymptote :
\(R^2\) hyperbolique ou linéaire selon disponibilité. \\
\textbf{Cohérence} & Corrélation Spearman entre fréquences temporelles
et spatiales, normalisée en {[}0,1{]}. \\
\textbf{Score pertinence} & Indicateur agrégé pondéré :
\(0.35 \\times Ir + 0.35 \\times \\text{complétude} + 0.20 \\times \\text{stabilité} + 0.10 \\times R^2\). \\
\textbf{AFC} & \textbf{Analyse Factorielle des Correspondances} :
méthode d'ordination exploratoire pour analyser les associations entre
catégories (ex. espèces et placettes). \\
\textbf{CA} & \textbf{Correspondence Analysis} : terme anglophone de
l'AFC. \\
\textbf{TSV} & \textbf{Tab-Separated Values} : format texte tabulaire
avec champs séparés par des tabulations. \\
\textbf{FAQ} & \textbf{Foire Aux Questions} : section de réponses aux
problèmes fréquents. \\
\textbf{R} & Langage et environnement de calcul statistique utilisé pour
ce pipeline. \\
\end{longtable}
}

\begin{center}\rule{0.5\linewidth}{0.5pt}\end{center}

\subsection{Références}\label{ruxe9fuxe9rences}

\begin{itemize}
\tightlist
\item
  Cours 24 Inventaires (diapositives N° 25-48 du ) - DU Mycologie 2026
  du Professeur Pierre-Arthur Moreau, Université de Lille (UFR3S PHAR)
\item
  Méthodes et définitions utilisées dans ce document (TEE, complétude,
  CA/AFC, cohérence tempo-spatiale) s'appuient sur les principes
  standards de l'écologie des communautés et de l'analyse multivariée
  sous R, en particulier via les packages \textbf{vegan}, \textbf{dplyr}
  et \textbf{ggplot2}. Pour un cadre théorique détaillé, se référer à la
  documentation officielle de ces packages ainsi qu'aux ouvrages de
  référence en ordination écologique.
\end{itemize}

\begin{center}\rule{0.5\linewidth}{0.5pt}\end{center}

\subsection{Contact}\label{contact}

\textbf{Auteur :} Eddy Boite\\
Pour les évolutions du script, ouvrir une issue ou documenter les
changements dans votre projet hôte.

\begin{center}\rule{0.5\linewidth}{0.5pt}\end{center}

\emph{Dernière mise à jour : 30 Mars 2026}
